\chapter*{Résumé}
\addcontentsline{toc}{chapter}{Résumé}

La modernisation des applications exige de comprendre l'existant, de résoudre les incompatibilités de versions, de transformer le code puis de revalider l'application après chaque évolution. Réalisé au sein de \textbf{CGI Technologies et Solutions Maroc}, ce projet de fin d'études vise à concevoir une \textbf{Migration Factory intelligente, extensible et gouvernée} afin d'automatiser une part importante de ce processus sans retirer au développeur le contrôle des décisions sensibles.

La contribution principale réside dans une \textbf{architecture SMA gouvernée}, hybride et centralement orchestrée, dans laquelle les opérations critiques restent sous contrôle déterministe et supervision humaine. Le projet a été conduit de manière incrémentale et itérative, depuis l'analyse du besoin et la conception de l'architecture jusqu'à l'implémentation, la vérification technique et l'exécution de scénarios représentatifs.

Les résultats obtenus montrent une trajectoire \textbf{Java/Spring Boot opérationnelle}, couvrant progressivement Spring Boot 2.1, 2.7, 3.5 et 4.0 avec les versions de Java adaptées aux stages. Trois démonstrations successives ont validé l'évolution du prototype en ligne de commande vers une \textit{Control Tower} full stack. La déclinaison Java intègre des espaces de travail isolés, des validations Maven, une boucle de réparation gouvernée, un cockpit de supervision, un assistant et un reporting assurant la conservation des preuves. Pour \textbf{Angular}, les transitions 18 vers 19 et 19 vers 20 ont été exécutées et scellées ; le stage 20 vers 21 est préparé mais non démarré. Ces résultats démontrent la capacité de la Factory à structurer une migration par stages, à relier les décisions aux preuves et à renforcer la traçabilité, permettant ainsi au développeur de superviser un processus plus explicite et reproductible.

Si la quantification des gains de temps, de coût et de productivité reste à consolider sur un corpus plus large, les résultats techniques obtenus démontrent déjà la faisabilité d'une migration gouvernée, traçable et supervisée sur les périmètres validés.

\vspace{0.4cm}
\noindent\textbf{Mots-clés :} Angular, Java/Spring Boot, migration logicielle, systèmes multi-agents, traçabilité.
